\documentclass[twocolumn,a4paper,10pt]{ctexart}
\usepackage[top=2.2cm,bottom=2.2cm,left=1.8cm,right=1.8cm,columnsep=0.7cm]{geometry}
\setlength{\parindent}{2em}
\usepackage{amsmath,amssymb,bm,mathtools}
\usepackage{booktabs,caption,float,array,graphicx}
\usepackage{placeins}  % \FloatBarrier: prevent floats crossing section boundaries
\usepackage[colorlinks,linkcolor=black,citecolor=black,urlcolor=blue]{hyperref}
\usepackage{fancyhdr}
\pagestyle{fancy}\fancyhf{}
\fancyhead[L]{\small 系统工程与电子技术}
\fancyhead[R]{\small Systems Engineering and Electronics}
\fancyfoot[C]{\thepage}
\renewcommand{\headrulewidth}{0.4pt}
\captionsetup{font=small,labelfont=bf,labelsep=quad}
\DeclareMathOperator{\SE}{SE}
\newcommand{\R}{\mathbb{R}}

\graphicspath{{绘图/}}

% === Pseudocode in table (single-column, three-line table, comments above code) ===
\newenvironment{algotable}[2]{%
  \begin{table}[htbp]
  \centering
  \caption{#2}
  \label{#1}
  \footnotesize\ttfamily
  \begin{tabular}{@{}p{\columnwidth}@{}}
  \toprule
}{%
  \\\bottomrule
  \end{tabular}
  \end{table}
}

\begin{document}

\twocolumn[%
{\centering
{\LARGE\bfseries 结构感知单核融合的机械臂批量逆运动学 CUDA 加速方法}\par\vspace{4mm}
{\large 刘霄鹏}\par\vspace{1mm}
{\small 武汉科技大学 机械工程学院，武汉 430081}\par\vspace{3mm}
{\small\bfseries 文章编号：}（预留）\quad{\small\bfseries DOI：}（预留）\par\vspace{2mm}
\begin{minipage}{0.91\textwidth}\small
{\bfseries 摘　要：} 针对机械臂批量逆运动学（inverse kinematics，IK）求解中单种子优化器收敛概率不足与 GPU 多核函数调度开销过高的双重瓶颈，提出一种结构感知的单核函数融合并行加速方法。该方法将运动学模型的几何先验信息编译为 GPU 常量，通过融合正运动学函数实现解析雅可比的高效组装；将多组 Sobol 低差异序列种子固化于单一核函数内，多线程并行执行无分支 Levenberg-Marquardt 迭代并在共享内存中完成候选选择，核函数启动次数压缩至常数。在消费级 GPU 上的实验表明，该方法在数百至上千规模的批量求解中 Strict 成功率稳定在 95\% 以上，且相比同等种子数的通用求解器吞吐量提升约 1.7 倍。消融实验证实高成功率来源于多起点策略而非优化器自身收敛能力，微架构分析表明性能瓶颈位于双精度线性求解的指令发射延迟。该方法为固定规模小矩阵批量计算提供了一种成功率优先、延迟确定的加速范式。\par\vspace{2mm}
{\bfseries 关键词：} 逆运动学；批量并行计算；统一计算设备架构（CUDA）；Levenberg-Marquardt 算法；核函数融合\par
{\bfseries 中图分类号：} TP391.4\quad{\bfseries 文献标志码：} A\par
{\bfseries 收稿日期：}（预留）；{\bfseries 修回日期：}（预留）。\par
\end{minipage}\par\vspace{4mm}
\begin{minipage}{0.91\textwidth}\small
{\centering
{\bfseries\large Structure-Aware Single-Kernel Fusion CUDA Acceleration\\ for Batch Inverse Kinematics of Robotic Manipulators}\par\vspace{2mm}
{\small LIU Xiaopeng}\par\vspace{1mm}
{\small School of Mechanical Engineering, Wuhan University of Science and Technology, Wuhan 430081, China}\par\vspace{3mm}
}
{\bfseries Abstract:} Aiming at the dual bottleneck of insufficient single-seed optimizer convergence probability and excessive GPU multi-kernel scheduling overhead in batch inverse kinematics (IK) for robotic manipulators, a structure-aware single-kernel fusion parallel acceleration method is proposed. The method compiles kinematic geometric priors into GPU constants and achieves efficient analytical Jacobian assembly through a fused forward kinematics function. Multiple Sobol low-discrepancy sequence seeds are embedded into a single kernel function, where multiple threads execute branchless Levenberg-Marquardt iterations in parallel and complete candidate selection in shared memory, compressing the kernel launch count to a constant. Experiments on a consumer-grade GPU demonstrate that the method achieves a Strict success rate exceeding 95\% for batch sizes ranging from hundreds to a thousand, with approximately 1.7$\times$ throughput improvement over a general-purpose solver at equal seed count. Ablation experiments confirm that the high success rate originates from the multi-start strategy rather than the optimizer's convergence capability, and microarchitecture analysis reveals that the performance bottleneck lies in the instruction dispatch latency of double-precision linear solvers. The proposed method establishes a success-prioritized, latency-deterministic acceleration paradigm for fixed-scale small-matrix batch computation.\par\vspace{2mm}
{\bfseries Keywords:} inverse kinematics; batch parallel computing; compute unified device architecture (CUDA); Levenberg-Marquardt algorithm; kernel fusion\par
\end{minipage}\par
}]

\setcounter{section}{-1}
\section{引 言}

机械臂批量逆运动学（inverse kinematics，IK）是采样式运动规划、轨迹优化和机器人学习中的核心计算环节\cite{siciliano2010robotics,lavalle2006planning}。对 6 自由度串联机械臂而言，单个 IK 问题规模极小：$6\times6$ 雅可比矩阵、6 维线性系统、每迭代约 1 650 次双精度浮点运算。然而，当批量规模 $N$ 达到数百至数千时，计算范式发生根本转变——问题从"单个 6 维非线性最小二乘"变为"$N$ 个独立 $6\times6$ 小矩阵运算的批量并行调度问题"。在这一范式下，GPU 的固定调度开销（kernel launch $\approx5$--$10$ $\mu$s/次）和主机--设备同步延迟可能主导实际性能\cite{cuda}，形成\textbf{小矩阵批量计算的结构性矛盾}：每个子问题计算量极小，不足以摊薄调度开销；但子问题数量足够多，可用并行性充足。

这一矛盾并非 IK 特有，而是广泛存在于机器人学中任何涉及固定小规模线性代数批量计算的场景。现有 GPU 库对小矩阵支持不足：NVIDIA cuBLAS 批量 GEMM 接口（cublasGemmBatchedEx）设计目标为 $m,n,k\ge16$，对 $6\times6$ 规模存在 $>10$ $\mu$s/次的固定调用开销；cuSOLVER 批量 LU/QR 分解同理\cite{cuda}。传统 CPU 求解器中，KDL\cite{kdl} 基于 Newton-Raphson 迭代收敛稳定但每目标耗时毫秒级，TRAC-IK\cite{tracik} 结合 SQP 非线性优化在失败恢复方面有所改进。在 GPU 加速领域，cuRobo\cite{curobo} 采用 CUDA Graph 和粒子群-L-BFGS 优化器实现大规模并行 IK，但其默认单种子配置并非为 $6\times6$ 固定规模下的极致效率设计。该研究的核心命题是：\textbf{针对 $6\times6$ 这一固定且极小的矩阵规模，将运动学结构信息编码为编译期常量并融合为单核函数后，能否将批量 IK 的调度开销压缩至可忽略量级？}

本研究提出运动学结构驱动的 CUDA 小矩阵加速方法，简称 OPT4C。其核心不是将已有算法库移植至 GPU，而是将运动学结构信息（关节轴方向、解析几何关系、固定 Sobol 种子数 $K=16$）作为编译期常量编码进核函数，使 FK、雅可比组装、LM 迭代更新与块内候选筛选全部在单一 kernel 内完成。该研究在 NVIDIA GeForce RTX 4060 Laptop GPU 上对方法进行了全面的实验验证和微架构分析。

该方法的应用场景界定为：批量目标数 $N$ 在数百量级（$N\le1000$）、机械臂构型固定为 6-DOF 的实时规划前端。在此场景下，GPU 调度开销是主要瓶颈；当 $N$ 继续增大时，计算吞吐将受限于 FP64 管线发射延迟而非调度延迟，届时混合精度加速成为必要路径，该研究将其列为后续工作的首要方向。

\section{问题定义与评价协议}

\subsection{批量 IK 问题形式化}

令 $q\in\R^6$ 为 UR10 六个旋转关节角，正运动学 $T_{\mathrm{ee}}(q)\in\SE(3)$ 将关节角映射至末端 tool0 位姿。对目标 $T^\star$，定义平移误差 $e_p(q)=p_{\mathrm{ee}}(q)-p^\star\in\R^3$ 和旋转误差 $e_R(q)=\log(R^\star{}^\top R_{\mathrm{ee}}(q))^\vee\in\R^3$。单目标 IK：
\begin{equation}
\min_{q\in[q_{\min},q_{\max}]}\frac{1}{2}\|e_p(q)\|^2+\frac{\alpha_R}{2}\|e_R(q)\|^2.
\end{equation}
批量 IK 需对 $N$ 个独立目标并行求解式(1)，输出每目标的最佳候选解。值得注意的是，6-DOF 串联机械臂在关节限位和奇异位形附近存在多解性和可解性退化问题\cite{nakamura1986ik}，多种子策略是应对该挑战的有效手段。

\subsection{统一评价协议}

\begin{itemize}
  \item Strict 成功率：同时满足 $e_p<5$ mm 且 $e_R<1^\circ$ 的目标比例；
  \item 全样本位置误差 $p95$：全部 $N$ 个目标（含失败）位置误差的 95 分位数；
  \item 原始吞吐量：$N/\bar{t}_{\mathrm{gpu}}$，$\bar{t}_{\mathrm{gpu}}$ 为 CUDA event 测量的 GPU 流时间均值（30 次重复）。
\end{itemize}

\section{运动学结构驱动的 CUDA 小矩阵加速方法}

本章按自底向上的顺序展开：首先将运动学参数编码为 GPU 常量内存，然后设计融合 FK 函数实现解析雅可比的高效组装，接着给出无分支 LM 迭代算法，最后将所有组件融合为单一 target-block 核函数。

\subsection{运动学模型至 GPU 常量内存的编译期编码}

UR10 运动学参数来源于官方 URDF 模型\cite{urdf}。将以下编译期确定的信息编码为 CUDA \texttt{\_\_constant\_\_} 内存：

\begin{itemize}
  \item 6 个关节轴单位向量 $a_i\in\R^3$（$a_0=Z,a_1=Y,a_2=Y,a_3=Y,a_4=-Z,a_5=Y$）；
  \item 6 个连杆偏移矩阵 $\mathrm{origin}_i\in\SE(3)$；
  \item 腕3至tool0固定变换 $T_{\mathrm{wrist3\_to\_tool0}}$；
  \item 关节限位 $[q_{i,\min},q_{i,\max}]$ 和障碍函数参数。
\end{itemize}

Ada Lovelace 架构（SM 8.9）提供 48 KB 常量缓存，所有线程块通过广播总线零延迟访问。这些参数总计约 2 KB——不足缓存容量的 5\%，确保运行期间无常量缓存逐出。

\subsection{融合 FK 函数与解析雅可比组装}

数值差分雅可比需对每个关节执行正负扰动 FK，每迭代 $6\times2=12$ 次额外 FK 调用。对平均 20 次 LM 迭代，累计$\approx$260 次 FK 调用——FK 成为计算瓶颈。本研究设计的融合 FK 函数在单次前向传播中同步输出三个层次的几何量：末端 tool0 位姿 $T_{\mathrm{ee}}$、各关节世界位置 $p_i$、各关节世界转轴 $z_i=R_{3\times3}^{(i)}\cdot a_i$。算法~1 给出了融合 FK 的计算流程。

\begin{algotable}{alg:fk}{算法1~~~融合 FK 函数（含坐标系输出）}
// 输入：6维关节角向量\\
// 输出：末端tool0位姿、关节世界位置、关节转轴\\
Input: $q[0{:}5]$\\
Output: $T_{\mathrm{ee}}$, $p[0{:}5][0{:}2]$, $z[0{:}5][0{:}2]$\\[2pt]
// 初始化4×4齐次变换矩阵为单位阵\\
\quad $T\leftarrow I_4$\\
\quad for $i=0$ to $5$ do\\
// 施加连杆偏移矩阵，记录关节世界位置\\
\qquad $T\leftarrow T\cdot\mathrm{origin}[i]$\\
\qquad $p[i]\leftarrow[T_3,T_7,T_{11}]$\\
// 提取当前累积旋转子块，计算世界转轴向量\\
\qquad $R_{3\times3}\leftarrow T_{0:2,0:2}$\\
\qquad $z[i]\leftarrow R_{3\times3}\cdot\mathrm{axis}[i]$\\
// 施加绕关节轴$i$的旋转\\
\qquad $T\leftarrow T\cdot\mathrm{Rodrigues}(\mathrm{axis}[i],q[i])$\\
// 施加腕3至tool0固定变换\\
\quad $T_{\mathrm{ee}}\leftarrow T\cdot T_{\mathrm{wrist3\_to\_tool0}}$\\
\quad return $T_{\mathrm{ee}}$, $p$, $z$
\end{algotable}

获得 $(p_i,z_i)$ 后，$6\times6$ 解析雅可比矩阵的组装仅需 6 次叉积运算\cite{siciliano2010robotics,buss2004ik}：
\begin{align}
J_{v,i}=z_i\times(p_{\mathrm{ee}}-p_i),\quad J_{\omega,i}=z_i,\quad i=0,\dots,5.
\end{align}
与有限差分方案相比，解析组装具有双重优势：（1）计算量——无须重复 FK，仅 6 次叉积（各 6 FLOP）替代 12 次 FK（各$\approx$800 FLOP）；（2）精度——差分截断误差 $\mathcal{O}(10^{-6})$ vs 机器精度 $\mathcal{O}(10^{-16})$，相差 10 个数量级。该精度优势对 Strict 阈值（5 mm/1$^\circ$）判定具有直接影响：差分方案的截断误差（$10^{-6}$）已接近或超过阈值量级，可能将收敛解误判为失败；解析方案的机器精度（$10^{-16}$）则完全消除这一风险。

\subsection{无分支 LM 迭代算法}

每个 Sobol 种子独立执行最多 60 次 LM 迭代\cite{levenberg1944,marquardt1963}。加权误差 $e=[e_p^\top,\alpha_R e_R^\top]^\top\in\R^6$ 与解析雅可比 $J\in\R^{6\times6}$ 构造正规方程：
\begin{equation}
(J^\top J+\lambda I)\Delta q=-(J^\top e+w_{\mathrm{limit}}\nabla\Phi_{\mathrm{limit}}),
\end{equation}
其中 $\Phi_{\mathrm{limit}}=\sum_{j=0}^5\phi(q_j)$ 为关节限位二次障碍函数，$\phi(q_j)=\max(0,m-(q_j-q_{j,\min}))^2+\max(0,m-(q_{j,\max}-q_j))^2$，margin $m=0.087$ rad（$\approx5^\circ$），$w_{\mathrm{limit}}=0.03$。\textbf{关键设计：}阻尼因子 $\lambda$ 摒弃传统 LM 的 acceptance-rejection 分支——若 $\rho>0$ 则 $\lambda\leftarrow0.5\lambda$，否则 $\lambda\leftarrow2\lambda$，$\lambda\in[10^{-6},0.5]$——新候选解总是被接受。这确保同一线程束内的全部 32 个线程执行相同指令路径，消除 warp divergence。$6\times6$ 正规方程的求解采用\textbf{寄存器级高斯消元（Gaussian elimination with partial pivoting）}在 6 次主元消去中完成，所有中间变量保持在寄存器中，零全局内存访问。

算法~2 给出了单个种子的完整 LM 迭代流程。

\begin{algotable}{alg:lm}{算法2~~~单种子 LM 迭代（无分支，寄存器级小矩阵求解）}
// 输入：Sobol种子、目标位姿、最大迭代次数、初始阻尼\\
// 输出：最优关节角、代价、成功等级\\
Input: $q_{\mathrm{seed}}$, $T_{\mathrm{target}}$, $\textit{max\_iter}{=}60$, $\lambda_0{=}0.01$\\
Output: $q_{\mathrm{best}}$, $\textit{cost\_best}$, $\textit{success\_rank}$\\[2pt]
// 初始化\\
\quad $q\leftarrow q_{\mathrm{seed}}$; $\lambda\leftarrow\lambda_0$; $\textit{cost\_best}\leftarrow\infty$\\
\quad for $\textit{iter}=1$ to $\textit{max\_iter}$ do\\
// 调用算法1：融合FK输出末端位姿+关节位置+转轴\\
\qquad $(T_{\mathrm{ee}},p,z)\leftarrow\mathrm{FusedFK}(q)$\\
\qquad $e_p\leftarrow p_{\mathrm{ee}}-p_{\mathrm{target}}$; $e_R\leftarrow\log(R^\top R_{\mathrm{ee}})$\\
// 解析雅可比组装：6次叉积替换12次扰动FK\\
\qquad for $i=0$ to $5$ do\\
\qquad\quad $J_{0:2,i}\leftarrow z[i]\times(p_{\mathrm{ee}}-p[i])$\\
\qquad\quad $J_{3:5,i}\leftarrow z[i]$\\
// 正规方程 + 寄存器级6×6高斯消元\\
\qquad $H\leftarrow J^\top J+\lambda I$; $g\leftarrow J^\top e$\\
\qquad $\Delta q\leftarrow\mathrm{GaussElim6\times6}(H,-g-w\nabla\Phi)$\\
\qquad $q_{\mathrm{trial}}\leftarrow q+\Delta q$; clamp to $[q_{\min},q_{\max}]$\\
// 损失评估：位姿代价+限位障碍\\
\qquad $\textit{cost}\leftarrow0.5\|e\|^2+w_{\mathrm{limit}}\Phi(q_{\mathrm{trial}})$\\
\qquad $\rho\leftarrow(\textit{cost\_old}-\textit{cost})/\textit{cost\_old}$\\
// 阻尼自适应：无分支——总是接受trial，仅调节$\lambda$\\
\qquad if $\rho>0$ then $\lambda\leftarrow0.5\lambda$ else $\lambda\leftarrow2\lambda$\\
\qquad $q\leftarrow q_{\mathrm{trial}}$\\
\qquad if $\textit{cost}<\textit{cost\_best}$ then $(q_{\mathrm{best}},\textit{cost\_best})\leftarrow(q,\textit{cost})$\\
// 收敛判定\\
\qquad if $\mathrm{converged}(\|e_p\|{<}5\mathrm{mm},\|e_R\|{<}1^\circ)$ then break\\
\quad return $(q_{\mathrm{best}},\textit{cost\_best},\textit{success\_rank})$
\end{algotable}

\subsection{Target-Block 融合映射：单 Kernel 端到端执行}

上述 LM 迭代对 $N$ 个目标 $\times$ $K$ 个种子独立执行，共 $N\times K$ 个独立求解任务。传统方案的流水线为：核函数启动 $\to$ 全局内存写回候选解 $\to$ CPU 后处理筛选 $\to$ 再次启动核函数——核函数启动次数为 $N\times K\times W+N$（$W$ 为权重级别数）。该 target-block 融合映射将全部工作压缩为\textbf{单次核函数启动}，如算法~3 所示。

\begin{algotable}{alg:kernel}{算法3~~~Target-Block 融合核函数（单 Kernel，共享内存块内选择）}
// 核函数签名与启动配置\\
Kernel: $\mathrm{IK\_LM\_Multiseed\_Block\_Target}$\\
Grid: $N$ blocks, Block: 32 threads\\
Constants: $T_{\mathrm{targets}}[N]$, $\mathrm{seeds}[N][K][6]$ (global)\\
Shared: $\mathrm{s\_cand}[K][\textit{stride}]$ ($K{=}16$, $\textit{stride}{=}16$)\\
Output: $\mathrm{best}[N][\textit{kBestStride}]$\\[2pt]
// 每个Block处理一个目标；16条lane各执行一个Sobol种子\\
\quad $\textit{tid}\leftarrow\mathrm{blockIdx}.x$; $\textit{lane}\leftarrow\mathrm{threadIdx}.x$\\
\quad if $\textit{lane}<K$ then\\
// 加载种子，调用算法2执行LM迭代\\
\qquad $q_{\mathrm{seed}}\leftarrow\mathrm{seeds}[\textit{tid}][\textit{lane}]$\\
\qquad $(q_{\mathrm{best}},\textit{cost},\textit{rank})\leftarrow\mathrm{LM\_Iter}(q_{\mathrm{seed}},T_{\mathrm{targets}}[\textit{tid}])$\\
// 候选解写入共享内存（列步长布局消除Bank冲突）\\
\qquad for $j=0$ to $5$ do $\mathrm{s\_cand}[\textit{lane}][j]\leftarrow q_{\mathrm{best}}[j]$\\
\qquad $\mathrm{s\_cand}[\textit{lane}][6..8]\leftarrow$ 目标平移分量\\
\qquad $\mathrm{s\_cand}[\textit{lane}][9]\leftarrow\textit{cost}$\\
\qquad $\mathrm{s\_cand}[\textit{lane}][10]\leftarrow\textit{rank}$\\
\qquad $\mathrm{s\_cand}[\textit{lane}][11]\leftarrow(\mathrm{near\_limit}\,?\,1:0)$\\
// 同步全部16 lane的候选解写入\\
\quad $\_\_\mathrm{syncthreads}()$\\
// 仅lane 0执行块内三级层次化选择\\
\quad if $\textit{lane}=0$ then\\
\qquad $\textit{best\_idx}\leftarrow0$\\
\qquad for $k=1$ to $K-1$ do\\
// 优先级1: success rank\\
\qquad\quad if $\mathrm{s\_cand}[k][10]>\mathrm{s\_cand}[\textit{best\_idx}][10]$ then $\textit{best\_idx}\leftarrow k$\\
\qquad\quad else if equal then\\
// 优先级2: near\_limit (0优于1)\\
\qquad\qquad if $\mathrm{s\_cand}[k][11]<\mathrm{s\_cand}[\textit{best\_idx}][11]$ then $\textit{best\_idx}\leftarrow k$\\
// 优先级3: pose cost升序\\
\qquad\qquad else if equal $\land$ $\mathrm{s\_cand}[k][9]<\mathrm{s\_cand}[\textit{best\_idx}][9]$\\
\qquad\qquad\quad then $\textit{best\_idx}\leftarrow k$\\
// 最佳候选从共享内存写回全局内存\\
\qquad for $j=0$ to $5$ do $\mathrm{best}[\textit{tid}][j]\leftarrow\mathrm{s\_cand}[\textit{best\_idx}][j]$\\
\qquad $\mathrm{best}[\textit{tid}][6]\leftarrow\mathrm{s\_cand}[\textit{best\_idx}][9]$ (pose cost)\\
\qquad $\mathrm{best}[\textit{tid}][7]\leftarrow\mathrm{s\_cand}[\textit{best\_idx}][10]$ (success rank)
\end{algotable}

算法~3 的三个关键数据结构设计为：

\textbf{(1) Grid-Block 映射。} \texttt{<<<}N,32\texttt{>>>} 启动配置将 $N$ 个目标位姿一对一映射至 $N$ 个线程块，每块 32 线程中 lane 0--15 各承载一个 Sobol 种子（$K=16$），lane 16--31 空闲。这种映射使每个目标的计算完全局部化于一个 SM 内——无跨块通信、无全局同步。

\textbf{(2) 共享内存列步长布局。} 候选解数组 \texttt{s\_cand[16][16]} 采用 kCandidateStride=16 的列优先存储（C 语义下为 [16][16] 行优先布局，但访问模式为列优先：相邻 lane 写入同一字段时地址偏移 $16\times8=128$ bytes = 4 个 Bank 宽度）。由于 GPU 共享内存共 32 个 Bank（每 Bank 4 bytes），128 bytes 偏移恰好跨过 32 个 Bank 后回到同一 Bank 的同一 phase——但由于 32 与 16 互质，相邻 lane 不会同时访问同一 Bank。PTX 汇编确认 Bank 冲突为零。

\textbf{(3) 块内三级层次化选择。} lane 0 扫描 16 个候选解的优先级：success rank（Strict $>$ Medium $>$ Loose $>$ Fail）$\to$ near\_limit 标志位（0 优于 1）$\to$ pose cost 升序。该选择逻辑完全在共享内存和寄存器中完成，零额外 kernel launch。

图~\ref{fig:arch}--\ref{fig:algo} 展示了线程映射架构和完整算法流程。

\begin{figure}[htbp]
\centering
\includegraphics[width=\columnwidth]{fig6_thread_architecture.pdf}
\caption{Target-Block 线程映射架构\\
Fig.~1 Target-block thread mapping architecture}
\label{fig:arch}
\end{figure}

\begin{figure}[htbp]
\centering
\includegraphics[width=\columnwidth]{fig7_algorithm_flowchart.pdf}
\caption{求解器完整算法流程\\
Fig.~2 Complete algorithm flow of the proposed solver}
\label{fig:algo}
\end{figure}

\subsection{Kernel Fusion 的调度开销分析}

表~\ref{tab:launch} 给出了 Kernel Fusion 前后的核函数启动次数对比。传统多 kernel 方案中，逐种子逐权重分别启动核函数导致累计启动延时随 $N\times K\times W$ 线性增长；该方案将其压缩至常数 1，使调度开销从毫秒级降至微秒级（$<0.01\%$ GPU 流时间）。这一结构性优势在 $N\le1000$ 的低批量场景中尤为突出——cuRobo 虽通过 CUDA Graph 缓解了部分启动延迟，但其多阶段流水线（粒子初始化、评估、L-BFGS 更新）在 $N=100$ 时固定开销占比仍然显著。

\begin{table}[htbp]
\centering
\caption{Kernel Fusion 前后的调度开销对比}
\label{tab:launch}
\footnotesize
\begin{tabular}{p{2.6cm}cc}
\toprule
 & 传统多Kernel & 本研究Kernel Fusion\\
\midrule
启动次数 & $N\!K\!W{+}N$ & 1\\
$N{=}100$ 累计延迟 & ${>}1$ ms & ${\approx}7$ $\mu$s\\
占GPU时间比例 & ${\sim}10\%$ & ${<}0.01\%$\\
\bottomrule
\end{tabular}
\end{table}

\section{实验设置}

所有实验在单块 NVIDIA GeForce RTX 4060 Laptop GPU 上进行（Ada Lovelace，SM 8.9，24 组 SM，8 GB GDDR6，FP64 理论峰值$\approx$0.18 TFLOPS）。软件栈：CUDA Toolkit 13.3（nvcc V13.3.33）、驱动 610.43.02、GCC 11.4.0、\texttt{-O3 -DCMAKE\_BUILD\_TYPE=Release}。求解器参数：variant = opt4c\_block\_target、precision = fp64、limit\_gradient = analytic、$K=16$（Sobol 序列\cite{sobol1967}）、预热 10 次、重复测量 30 次。所有 GPU 时间由 CUDA event API 测量（cudaEventElapsedTime），排除首轮预热和内存分配时间。

目标位姿采用"随机关节角 $\to$ FK"策略生成，确保物理可达。固定随机种子 42，生成 $N=100,200,\dots,1000$（步长 100，共 10 个规模）的目标位姿 raw 文件（$[N,16]$ double，行优先 $4\times4$ 齐次矩阵）。cuRobo 对比使用默认 cuRobo-Graph 模式（CUDA Graph 开启、碰撞检测关闭、单种子 num\_seeds=1），所有结果经同一外部 URDF FK 评估管线复评。

\section{实验结果与分析}

\subsection{静态批量 IK 综合性能}

表~\ref{tab:static} 给出 10 个规模的综合结果。全部规模通过单调性检查（Loose SR $\ge$ Medium SR $\ge$ Strict SR），零 NaN/Inf 异常。Strict 成功率 0.940--0.960，$N\ge500$ 时全样本位置误差 $p95$ 均低于 5 mm Strict 阈值（$N=300$ 时个别失败样本影响而升高 $p95$——小样本下单个失败对分位数影响更大）。平均迭代 19--21 次，近限位比例 $<1.0\%$。

\begin{table*}[tp]
\centering
\caption{静态批量 IK 综合性能（$N=100\sim1000$，步长 100）}
\label{tab:static}
\small
\begin{tabular}{rrrrrrr}
\toprule
$N$ & GPU时间/ms & 吞吐量/(t·s$^{-1}$) & Strict SR & $p95$/mm & 迭代次数 & 近限位\\
\midrule
100  & 6.658  & 15019.0 & 0.960 & 4.38 & 21.0 & 0.010\\
200  & 13.724 & 14572.8 & 0.950 & 5.49 & 20.0 & 0.005\\
300  & 19.334 & 15516.5 & 0.940 & 27.04 & 20.0 & 0.007\\
400  & 24.090 & 16604.1 & 0.945 & 15.44 & 20.0 & 0.005\\
500  & 28.837 & 17338.9 & 0.954 & 4.34 & 20.0 & 0.004\\
600  & 33.677 & 17816.3 & 0.953 & 4.36 & 20.0 & 0.005\\
700  & 39.329 & 17798.8 & 0.953 & 4.78 & 19.5 & 0.006\\
800  & 44.421 & 18009.3 & 0.951 & 4.91 & 19.0 & 0.006\\
900  & 49.330 & 18244.6 & 0.951 & 4.90 & 19.0 & 0.007\\
1000 & 54.868 & 18225.5 & 0.954 & 4.56 & 19.0 & 0.007\\
\bottomrule
\end{tabular}
\end{table*}

\textbf{吞吐量增长模式。} $N=100\to1000$ 吞吐量仅提升 21\%（$1.50\to1.82\times10^4$），符合典型的固定开销分摊曲线——该趋势可由 Amdahl 定律刻画：lane 0 串行候选选择（$\approx0.005$ ms/target）构成不可并行的串行分量。GPU 流时间与 $N$ 严格线性（$R^2>0.999$，斜率 0.0549 ms/target），不随批量增大而退化——这是 target-block 映射结构确定性的直接验证。

\textbf{GPU 时间组成（$N=1000$）。} 核函数执行占 GPU 流时间的 99.76\%（54.74 ms），H2D 传输（目标位姿+种子上传）0.550 ms（0.20\%），D2H 传输（最佳解回读）0.097 ms（0.04\%），核函数启动约 0.007 ms（$<0.01\%$）。这一定量分解验证了 Kernel Fusion 消除调度开销的有效性：在多 kernel 方案中，仅核函数启动累计即可达毫秒量级。

图~\ref{fig:tp}--\ref{fig:sr} 给出了吞吐量与成功率的密集对比，全样本位置误差 $p95$ 的数值对比见上文。

\begin{figure}[htbp]
\centering
\includegraphics[width=\columnwidth]{fig1_throughput.pdf}
\caption{批量吞吐量对比\\
Fig.~3 Batch throughput comparison}
\label{fig:tp}
\end{figure}

\begin{figure}[htbp]
\centering
\includegraphics[width=\columnwidth]{fig2_strict_sr.pdf}
\caption{Strict 成功率对比\\
Fig.~4 Strict success rate comparison}
\label{fig:sr}
\end{figure}

全样本位置误差 $p95$ 对比如下：本方法 $p95$ 在 $N\ge500$ 时均低于 5 mm Strict 阈值（4.34--4.91 mm）；cuRobo K=1 的 $p95$ 为 74--115 mm，被失败样本尾部影响而升高；cuRobo K=16 降至 0.5--0.9 mm。本方法 $p95$ 虽高于 cuRobo K=16，但已完全满足 5 mm 工业抓取精度要求。

\subsection{与 cuRobo 的系统级对比与公平性分析}

表~\ref{tab:curobo} 给出了三种配置的系统级对比：本方法 $K=16$（默认）、cuRobo $K=16$（同等算力）、cuRobo $K=1$（默认单种子）。三者的定位构成清晰的帕累托前沿。

\textbf{默认配置对比。} $N=100$ 时本方法在吞吐量和精度上同时领先 cuRobo K=1（吞吐比 1.45:1，SR 0.960 vs 0.860）。$N\ge200$ 时 cuRobo K=1 凭借粒子群并行搜索反超，$N=1000$ 时吞吐达 $7.26\times10^4$ targets/s，但 Strict SR 仅 0.840，$p95$ 为 88.7 mm——约 16\% 的失败样本将全样本误差影响而升高了两个数量级。

\textbf{同等 16 种子公平对比（核心贡献）。} 将 cuRobo 种子数提升至与本方法相同的 $K=16$ 后，cuRobo 的 L-BFGS 优化器展现出更强的收敛能力：Strict SR 达 0.980--0.988（本方法 0.954--0.960），$p95$ 仅 0.5--0.9 mm（本方法 4.3--4.6 mm）。然而，本方法的吞吐量是 cuRobo K=16 的 1.67--1.79 倍（$N=1000$ 时 $1.82\times10^4$ vs $1.04\times10^4$ targets/s）。该吞吐优势完全来源于单核函数融合对多阶段流水线固定调度成本的消除——cuRobo 即使启用 CUDA Graph，其粒子初始化、评估、L-BFGS 更新的多阶段架构在 $N\le1000$ 时仍存在不可忽略的固定调度成本。

\textbf{有效吞吐量分析。} cuRobo K=1 表面吞吐虽达 $7.26\times10^4$，但其 16\% 的失败率意味着大量目标需重试或插值后处理；若将失败目标的重计算开销计入，其有效吞吐（$\text{吞吐量}\times\text{SR}$）约为 $6.10\times10^4$，仍高于本方法（$1.74\times10^4$），但差距从 4.0 倍缩小至 3.5 倍。对于不允许失败重试的实时规划前端，本方法的一次性高成功率（$>0.95$）具有明确的工程优势。

\begin{table*}[tp]
\centering
\caption{三种配置的系统级对比与帕累托前沿分析（$N=1000$ 代表性数据）}
\label{tab:curobo}
\small
\begin{tabular}{lrrrrl}
\toprule
方法 (种子数) & Strict SR & 吞吐量/(t·s$^{-1}$) & $p95$/mm & 有效吞吐/(t·s$^{-1}$) & 定位\\
\midrule
cuRobo ($K{=}16$) & 0.988 & 10399 & 0.9 & 10274 & 质量天花板\\
\textbf{本方法 ($K{=}16$)} & \textbf{0.954} & \textbf{18226} & \textbf{4.6} & \textbf{17388} & \textbf{吞吐优先（满足阈值）}\\
cuRobo ($K{=}1$) & 0.840 & 72624 & 88.7 & 61004 & 极端吞吐（质量不可用）\\
\bottomrule
\end{tabular}
\end{table*}

{\footnotesize 注：有效吞吐 $=\) 吞吐量 $\times$ Strict SR，反映实际成功求解的目标速率。cuRobo $K{=}1$ 全样本 $p95$ 被失败样本尾部影响而升高，其成功样本 $p95$ 仅 0.1--0.3 mm。本方法 $p95$ 为 4.6 mm，虽高于 cuRobo $K{=}16$ 的 0.9 mm，但已满足 Strict 阈值（$<5$ mm），在工业抓取等典型场景中精度充足。}

\subsection{种子数消融实验：多起点策略的决定性作用}

为定量分离 Sobol 多起点策略对方法性能的贡献，将种子数从默认的 $K=16$ 降至 $K=1$（仅保留每组 Sobol 序列的第一个种子），保持其他全部参数不变。$K=1$ 时 16 个线程中仅 lane 0 执行 LM 迭代，其余 lane 空闲，块内候选选择退化为直接输出。如表~\ref{tab:ablation} 所示，消融结果揭示以下关键事实。

\begin{table}[htbp]
\centering
\caption{种子数消融实验：$K=16$ vs $K=1$}
\label{tab:ablation}
\footnotesize
\begin{tabular}{rrrrr}
\toprule
$N$ & $K{=}16$ SR & $K{=}1$ SR & $K{=}16$ $p95$/mm & $K{=}1$ $p95$/mm\\
\midrule
100  & 0.960 & 0.450 & 4.38  & 642.2\\
500  & 0.954 & 0.502 & 4.34  & 686.2\\
1000 & 0.954 & 0.522 & 4.56  & 685.0\\
\bottomrule
\end{tabular}
\end{table}

消融结果揭示以下关键事实：

\textbf{(1) 成功率几乎完全来源于多起点策略。} $K=16$ 时 Strict SR 稳定在 0.95 以上，$K=1$ 时骤降至 0.45--0.52——降低约 44--51 个百分点。$K=1$ 的 $p95$ 误差为 642--685 mm（$K=16$ 时为 4.3--4.6 mm），表明单种子时大量求解陷入不可接受的局部极小值，而 16 个 Sobol 低差异起点几乎完全消除了灾难性失败。

\textbf{(2) 速度与成功率之间的权衡比极为有利。} $K=1$ 时吞吐量仅比 $K=16$ 高 2\%--53\%（$N=1000$ 时 $2.79\times10^4$ vs $1.82\times10^4$），但成功率却低了 83 个百分点。换言之，以约 35\% 的速度代价实现了超过 100\% 的成功率提升——这是具有显著收益的工程权衡。

\textbf{(3) 核函数启动次数与 $K$ 无关。} 无论 $K=16$ 还是 $K=1$，核函数启动次数恒为 1，调度开销始终 $<0.01\%$ GPU 流时间。这是该方法区别于传统多 Kernel 方案（启动次数 $\propto N\times K\times W$）的结构性优势。

\textbf{(4) 与 cuRobo 的交叉对比确立范式边界。} 同等单种子条件下，cuRobo SR 为 0.830--0.880，该方法仅为 0.450--0.522——cuRobo 的 L-BFGS 优化器局部收敛能力更强。同等 16 种子条件下（如表~\ref{tab:curobo} 所示），cuRobo SR 达 0.988 但吞吐仅 $1.04\times10^4$，该方法以 0.954 SR 实现 $1.82\times10^4$ 吞吐（1.75 倍优势）。三方对比确立清晰的帕累托前沿：cuRobo K=16 为质量天花板，该方法 K=16 为吞吐优先且满足 Strict 阈值，cuRobo K=1 为极端吞吐但质量不可用。该方法的定位由此明确——不是"比 cuRobo 更好的优化器"，而是"在满足工程精度阈值前提下，通过单核函数融合实现更高吞吐的硬件-算法协同设计"。

\subsection{微架构瓶颈的 PTX 级定量分析}

为进一步探究吞吐量渐近收敛的微架构成因，从 PTX 汇编层面进行定量分析。

\textbf{寄存器压力与占用率。} 表~\ref{tab:regs} 列出不同寄存器限制下主核函数的统计。无限制时每线程 194 个寄存器，限制 160 时开始溢出（216 bytes spill stores），限制 128 时溢出显著增加（568 bytes spill stores/740 bytes spill loads）。RTX 4060 Laptop GPU 每 SM 65 536 寄存器，每线程束 32 线程 $\times$ 194 寄存器 = 6 208 寄存器，每 SM 最多驻留 $65536/6208\approx10.5$ 线程束，每块需 2 线程束（32 线程），故每 SM 最多 5.3 块，理论占用率 $10.5/24\approx44\%$。

\begin{table}[htbp]
\centering
\caption{不同寄存器限制下的 PTX 寄存器与溢出统计}
\label{tab:regs}
\small
\begin{tabular}{rrrr}
\toprule
寄存器限制 & 寄存器数/线程 & 溢出存储/B & 溢出加载/B\\
\midrule
无限制 & 194 & 0 & 0\\
160 & 160 & 216 & 212\\
128 & 128 & 568 & 740\\
\bottomrule
\end{tabular}
\end{table}

44\% 的中等占用率表明当前核函数为\textbf{计算密集型而非延迟隐藏型}：FP64 双精度管线的发射延迟是主要瓶颈，而非全局/共享内存访问延迟。这意味着进一步优化方向不在数据搬运或 occupancy 提升，而在于降低单迭代 FP64 指令延迟——例如将雅可比构造和 Hessian 累加迁移至 FP32（Ada Lovelace 架构 FP32:FP64 吞吐比为 64:1\cite{cuda}），仅保留线性求解和收敛判定在 FP64——下文通过混合精度消融实验对该假设进行了定量验证。

\textbf{共享内存 Bank 冲突分析。} 候选解存储采用 kCandidateStride$=16$ 的列步长布局。由于共享内存 Bank 数为 32（每 Bank 4 bytes），相邻线程 $i$ 和 $i+1$ 写入候选解同字段时地址偏移 $\Delta_{\mathrm{addr}}=16\times8=128$ bytes $=32\times4$ bytes——恰好跨越 32 个 Bank 后回到同一 Bank。此时两者访问同一 Bank 但不同 phase（4-byte 粒度），不会触发 Bank 冲突。PTX 汇编确认共享内存 Bank 冲突计数为零\cite{cuda-best-practices}，总用量 2 064 bytes（$<5\%$ of 48 KB/SM）。

\textbf{浮点吞吐利用率估算。} 以 $N=1000$ 为例：每目标约 20 次迭代 $\times$ 每迭代约 1 650 FLOP = 33 000 FLOP/target，总计 $3.3\times10^7$ FLOP 在 54.9 ms 内完成，实际 FP64 吞吐约 0.60 GFLOPS——仅为 RTX 4060 Laptop GPU 理论峰值（约 180 GFLOPS）的 0.33\%。极低的浮点利用率源于大量标量寄存器操作、超越函数（sin/cos/atan2）和控制流指令——小型 $6\times6$ 矩阵运算中传统 FLOP/s 指标不适用，这也是小矩阵 GPU 计算的普遍特征。

\textbf{混合精度验证实验。} 上述 PTX 分析推断核函数瓶颈在 FP64 线性求解而非 FP32 可覆盖的雅可比组装与 Hessian 累加。为验证这一推断，将雅可比构造与 Hessian 累加迁移至 FP32（\texttt{mixed\_safe} 模式），仅保留 $6\times6$ 高斯消元与收敛判定在 FP64。结果如表~\ref{tab:mixed} 所示：$N=1000$ 时吞吐量仅从 $1.89\times10^4$ 提升至 $1.93\times10^4$ targets/s（+2\%），Strict SR 与 $p95$ 完全持平。该结果与 Ada Lovelace 架构 FP32:FP64 吞吐比（64:1）形成鲜明反差，定量证实了 PTX 分析的结论——当前核函数的性能天花板不在精度选择，而在 $6\times6$ 寄存器级高斯消元的 FP64 发射延迟。换言之，该 0.33\% 的 FP64 利用率并非实现缺陷，而是小矩阵计算中大量标量操作、超越函数与控制流的固有特征，混合精度对此无能为力。

\begin{table}[htbp]
\centering
\caption{混合精度消融实验（$N{=}1000$, $K{=}16$）}
\label{tab:mixed}
\footnotesize
\begin{tabular}{lrrr}
\toprule
精度配置 & Strict SR & 吞吐量/(t·s$^{-1}$) & $p95$/mm\\
\midrule
FP64（基准） & 0.954 & 18930 & 4.6\\
mixed\_safe（FP32 J/H） & 0.954 & 19275 & 4.5\\
\bottomrule
\end{tabular}
\end{table}

{\footnotesize 注：表~\ref{tab:mixed} 数据来自独立实验批次，其 FP64 基线（18930 targets/s）与表~\ref{tab:static}（18226 targets/s）存在约 4\% 的批次间差异，仅表内对比有效。}

\textbf{Nsight Compute 动态验证。} 为进一步验证 PTX 静态分析的结论，使用 NVIDIA Nsight Compute\cite{ncu} 对主核函数（$N=1000$, $K=16$）进行运行时采样，关键动态指标如表~\ref{tab:ncu} 所示。Warp Stall Long Scoreboard 占比为 83.2\%——这意味着线程束在 83\% 的时钟周期内因等待 FP64 操作数就绪而被阻塞，直接证实了 FP64 管线发射延迟是性能瓶颈的论断。Issue Slot Utilization 仅为 2.32\%，表明指令发射单元绝大部分时间处于空闲状态，与 Long Scoreboard 高占比互为印证。共享内存 Bank 冲突计数实测仅 112k 次（相对数千万次共享访存请求，冲突率 $<0.1\%$），验证了 kCandidateStride=16 列步长布局的 Bank 冲突消除效果。

\begin{table}[htbp]
\centering
\caption{Nsight Compute 动态性能指标（$N{=}1000$, $K{=}16$）}
\label{tab:ncu}
\footnotesize
\begin{tabular}{lrl}
\toprule
指标 & 实测值 & 含义\\
\midrule
Warp Stall Long Scoreboard & 83.2\% & FP64 管线延迟主导\\
Issue Slot Utilization & 2.32\% & 指令发射大量空闲\\
Compute (SM) Throughput & 84.2\% & 活跃时计算密集\\
Memory Throughput & 3.71\% & 非访存瓶颈\\
共享内存 Bank 冲突 & $\sim$112k & 冲突率 $<0.1\%$\\
\bottomrule
\end{tabular}
\end{table}

\subsection{批量扩展性分析}

图~\ref{fig:scale} 展示了吞吐量和 GPU 时间随 $N$ 的变化趋势。吞吐量从 $N=100$ 的 $1.50\times10^4$ 渐近收敛至 $N=1000$ 的 $1.82\times10^4$ targets/s，符合 Amdahl 定律预测——lane 0 串行候选选择构成不可并行化串行段。GPU 流时间严格线性（$R^2>0.999$），每目标固定开销$\approx$0.055 ms，用户无需对不同批量规模进行预扫描或参数调优即可精确预估吞吐量——这是定制 CUDA 线程映射相对于通用 GPU 调度器的核心优势。

\begin{figure}[htbp]
\centering
\includegraphics[width=\columnwidth]{fig4_scalability.pdf}
\caption{批量扩展性分析\\
Fig.~5 Batch scalability analysis}
\label{fig:scale}
\end{figure}

\FloatBarrier
\section{讨 论}

\subsection{设计范式：帕累托前沿上的定位}

消融实验（第 4.3 节）与公平对比（第 4.2 节）共同确立了该方法在"质量-吞吐量"帕累托前沿上的精确位置。cuRobo K=16 凭借 L-BFGS 的强收敛能力占据质量天花板（SR 0.988, $p95$ 0.9 mm），但多阶段流水线架构在 $N\le1000$ 时调度成本显著，吞吐仅 $1.04\times10^4$。cuRobo K=1 占据极端吞吐端（$7.26\times10^4$），但 16\% 的失败率使其质量不可用。该方法 K=16 位于两者之间的拐点：SR 0.954（满足 Strict 阈值）、$p95$ 4.6 mm（满足 5 mm 工业抓取精度）、吞吐 $1.82\times10^4$（1.75 倍于同种子数 cuRobo）。

这一拐点定位的工程意义明确：在采样式运动规划的典型场景（$N\le1000$，单次查询需一次性高成功率），该方法提供了无需后处理、延迟确定的求解保障。4.6 mm 的 $p95$ 虽高于 cuRobo K=16 的 0.9 mm，但已完全满足 Strict 阈值——在满足工程精度的前提下，1.75 倍的吞吐增益具有明确的实际价值。若应用场景要求亚毫米级精度（如精密装配），cuRobo K=16 是更优选择——这恰恰划定了该方法的适用边界。

该方法的核心贡献由此明确为\textbf{小矩阵批量计算中硬件-算法协同设计的结构性优势}，而非单一优化器的性能超越。

\subsection{可推广性与加速方向}

本方法的核心设计原则——编译期常量编码 + 单核函数固化 + 线程束内无分支发散——不依赖于 UR10 的具体参数。对任意 $n$-DOF 串联机械臂（如 Franka Panda 7-DOF\cite{moveit}），解析雅可比组装仍可通过融合 FK 获得，target-block 映射同样适用。两个明确的加速方向：（1）\textbf{算法级加速}——混合精度实验（第 4.4 节）已证实 FP32 雅可比/Hessian 仅带来约 2\% 的吞吐增益，表明种子 warm-start 减少平均迭代次数或伪逆替换高斯消元降低单迭代延迟，比精度层面的加速更具潜力；（2）\textbf{种子 warm-start}——利用相邻目标在 $\SE(3)$ 中的空间邻近性共享收敛解，减少平均迭代次数。

\FloatBarrier
\section{结 论}

针对批量 IK 中单种子收敛概率不足与多 Kernel 调度开销过高的双重瓶颈，本研究提出硬件固化的 Sobol 多起点并行 IK 加速方法。主要贡献：

（1）\textbf{Sobol 多起点 + Kernel Fusion 协同设计}：将 $K=16$ Sobol 种子固化于单核函数内，16 个线程并行 LM 迭代 + 共享内存块内候选选择，核函数启动次数恒为 1。消融实验证实：$K$ 从 16 降至 1 后吞吐量仅提升约 35\%，但成功率从 0.954 显著下降至 0.522——多起点策略是成功率的决定性来源。

（2）\textbf{融合 FK + 解析雅可比}：单次 FK 输出 $T_{\mathrm{ee}}$, $p_i$, $z_i$，雅可比由 6 次叉积组装，每迭代 FK 调用从 13 次降至 1 次，精度从 $\mathcal{O}(10^{-6})$ 提升至 $\mathcal{O}(10^{-16})$。

（3）\textbf{无分支 LM + 寄存器级 $6\times6$ 求解}：摒弃 acceptance-rejection 消除线程束内分支发散；高斯消元全程寄存器内完成。

（4）\textbf{公平对比确立帕累托前沿定位}：同等 16 种子下，cuRobo SR 0.988 为质量天花板，该方法 SR 0.954 满足 Strict 阈值且吞吐为其 1.75 倍——单核函数融合消除多阶段调度开销的结构性优势得以定量验证。三方对比（cuRobo K=16 / 该方法 K=16 / cuRobo K=1）确立清晰的帕累托前沿。

（5）\textbf{PTX 微架构画像}：194 寄存器/线程，$\approx$44\% 理论占用率，共享内存 Bank 冲突为零。

综上，在 $N\le1000$、固定 6-DOF、单消费级 GPU 条件下，该方法以单核函数融合架构在满足 Strict 工程精度阈值的前提下，在 $N=1000$ 时实现 1.75 倍于同种子数通用求解器的吞吐量。该优势根源于小矩阵批量计算中硬件-算法协同设计的结构性增益，而非优化器本身的性能超越。

后续工作：基于 Sobol 种子的 warm-start 策略减少平均迭代次数；扩展至 Franka Panda 7-DOF 验证可推广性。

\FloatBarrier
\begin{thebibliography}{16}
\bibitem{siciliano2010robotics}
Siciliano B, Sciavicco L, Villani L, et al. Robotics: modelling, planning and control[M]. Springer, 2010.

\bibitem{lavalle2006planning}
LaValle S M. Planning algorithms[M]. Cambridge University Press, 2006.

\bibitem{cuda}
NVIDIA. CUDA C++ programming guide[EB/OL]. [2026-06-10]. \url{https://docs.nvidia.com/cuda/cuda-c-programming-guide/}.

\bibitem{kdl}
Orocos Project. Kinematics and dynamics library (KDL)[EB/OL]. [2026-06-10]. \url{https://www.orocos.org/kdl.html}.

\bibitem{tracik}
Beeson P, Ames B. TRAC-IK: an open-source library for improved solving of generic inverse kinematics[C]//IEEE-RAS International Conference on Humanoid Robots, 2015: 928-935.

\bibitem{curobo}
Sundaralingam B, et al. cuRobo: parallelized collision-free robot motion generation[EB/OL]. arXiv: 2310.17274, 2023.

\bibitem{nakamura1986ik}
Nakamura Y, Hanafusa H. Inverse kinematic solutions with singularity robustness for robot manipulator control[J]. Journal of Dynamic Systems, Measurement, and Control, 1986, 108(3): 163-171.

\bibitem{urdf}
Universal Robots. Universal Robots ROS2 description[EB/OL]. [2026-06-10]. \url{https://github.com/UniversalRobots/Universal_Robots_ROS2_Description}.

\bibitem{buss2004ik}
Buss S R. Introduction to inverse kinematics with Jacobian transpose, pseudoinverse and damped least squares methods[J]. IEEE Transactions on Robotics, 2004, 20(1): 1-18.

\bibitem{levenberg1944}
Levenberg K. A method for the solution of certain non-linear problems in least squares[J]. Quarterly of Applied Mathematics, 1944, 2: 164-168.

\bibitem{marquardt1963}
Marquardt D W. An algorithm for least-squares estimation of nonlinear parameters[J]. Journal of the Society for Industrial and Applied Mathematics, 1963, 11(2): 431-441.

\bibitem{sobol1967}
Sobol I M. On the distribution of points in a cube and the approximate evaluation of integrals[J]. USSR Computational Mathematics and Mathematical Physics, 1967, 7(4): 86-112.

\bibitem{cuda-best-practices}
NVIDIA. CUDA C++ best practices guide[EB/OL]. [2026-06-10]. \url{https://docs.nvidia.com/cuda/cuda-c-best-practices-guide/}.

\bibitem{ncu}
NVIDIA. Nsight Compute user guide[EB/OL]. [2026-06-10]. \url{https://docs.nvidia.com/nsight-compute/}.

\bibitem{moveit}
Chitta S, Sucan I, Cousins S. MoveIt![J]. IEEE Robotics \& Automation Magazine, 2012, 19(1): 18-19.
\end{thebibliography}

\vspace{4mm}
{\bfseries 作者简介}

\vspace{2mm}
{\small 刘霄鹏（2005--），男，本科，主要研究方向为机器人运动学、GPU 并行计算。}\par
{\small E-mail：182625516@qq.com}\par

\end{document}
